% Owner: Huy (Phase 1). Lab item 3.1; docs/application/PROJECT.md §2.1 and §4.
\section{Introduction}
\label{sec:intro}

Academic regulations at Vietnamese universities cover training, academic
warnings, graduation and similar matters. They are issued as long documents
divided into articles (\emph{Điều}) and clauses (\emph{Khoản}). Students ask
about these rules in everyday language, so a question is often worded quite
differently from the article that answers it. This mismatch is known as the
\emph{lexical gap}. A literal search with Ctrl+F in a PDF lists matching
strings in document order. A student searching for a phrase that the relevant
article never uses may therefore miss it.
General-purpose chatbots answer fluently, but they may produce statements that
are not in the regulations or cite an article that does not support their answer.

% Example: data/ViRHE4QA_v2.csv of https://github.com/longstnguyen/ViRE at commit
% 60ea7155134eaa426ab1d0a82d2bbbf327d1b470 (file sha256 23fb969985150599d3ad0b5229d039b8541ed335fbc29dcfa0b0d733ccaa5682),
% row index=9542 (original split: train). Checks and their outputs:
% data/raw/virhe4qa/checks/ (local, git-ignored). Quotes keep the dataset wording;
% only whitespace is normalised. Its context group must go to D1-dev in our
% split, because it is used here (docs/application/PROJECT.md, Section 6.2).
For example, one question in the public ViRHE4QA
subset~\cite{do2024r2gqaretrieverreadergeneratorquestionanswering,nguyen-quan-2026-works}
asks ``Việc làm KLTN có thể bị dừng giữa chừng không?'', that is, whether work
on a graduation thesis can be stopped midway (KLTN abbreviates \emph{khóa luận
tốt nghiệp}, graduation thesis). The answer is in an article titled ``Báo cáo
tiến độ thực hiện KLTN'' (progress reports on the thesis), which lists ``Đình
chỉ việc thực hiện KLTN'' (suspending the thesis) among the possible measures. The
student's phrase \emph{dừng giữa chừng} (stop midway) does not occur in the
article, which uses the legal term \emph{đình chỉ} (suspend) instead, and the
title of the article does not mention stopping at all.

We cast the task as \emph{semantic retrieval}. Given a Vietnamese question $q$
and a collection of regulation passages, the system ranks the articles by
relevance and returns the top few, each with the text of its best-matching
clause and its source document.
Machine learning suits this problem because the set of possible paraphrases of
a rule is open-ended. Hand-written synonym lists cannot cover it, whereas a
learned sentence encoder maps texts with similar meaning to nearby vectors.
The learning setting has two levels. The encoder is trained by
\emph{contrastive representation learning}: self-supervised on raw sentences in
unsupervised SimCSE, and supervised by labelled sentence pairs in supervised
SimCSE (Section~\ref{sec:method}). Our application does not train the encoder
further. It uses a pretrained encoder zero-shot to score question--passage
pairs and combines these scores with BM25 lexical matching~\cite{Robertson_2009}.
The only quantities we fit are a score-fusion weight and a confidence
threshold. Both are selected on a development split of a labelled retrieval
dataset, and the test data are used only for the final report
(Section~\ref{sec:setup}).

\textbf{Selected paper.} We study \emph{SimCSE: Simple Contrastive Learning of
Sentence Embeddings} by Tianyu Gao, Xingcheng Yao and Danqi Chen, published in
the Proceedings of the 2021 Conference on Empirical Methods in Natural Language
Processing (EMNLP 2021), pages 6894--6910,
DOI~\url{10.18653/v1/2021.emnlp-main.552}~\cite{gao-etal-2021-simcse}.
The CORE portal ranks EMNLP as A in CORE2021 and as A* in CORE2023 and
ICORE2026.\footnote{\url{https://portal.core.edu.au/conf-ranks/448/}}
We chose this paper for four reasons.
\begin{enumerate}
\item It learns exactly the representation that our retrieval needs: sentence
embeddings whose cosine similarity reflects meaning.
\item The method is simple. It uses one loss, one encoder and standard dropout
as the only augmentation, so it can be analysed in full and re-implemented
within a course project.
\item It reports large gains. With BERT-base, the average Spearman correlation
over seven semantic textual similarity (STS) tasks is 76.25 for unsupervised
and 81.57 for supervised SimCSE, against 56.70 for averaged BERT embeddings and
74.89 for SBERT~\cite{reimers-gurevych-2019-sentence}, as reported in Table~5
of~\cite{gao-etal-2021-simcse}.
\item Its code and checkpoints are public, so the main results can be
re-evaluated on a laptop CPU. A community Vietnamese checkpoint, trained
following the SimCSE recipe on PhoBERT~\cite{nguyen-tuan-nguyen-2020-phobert},
is also available for the application.
\end{enumerate}

\textbf{Contributions.} The parts reused from earlier work and the parts we
built are listed separately in Section~\ref{sec:app}. This report
\begin{itemize}
\item analyses SimCSE: its training objective, the role of dropout noise and
hard negatives, its link to alignment, uniformity and anisotropy, and the
execution flow of the official code at commit \texttt{235f89e}
(Section~\ref{sec:method});
% Numbers below: experiments/reproduction/results/R{1,2,3}/metrics.json ("Avg.").
\item reproduces the main STS results with the official evaluation script and
released checkpoints. The averages we obtain are 81.58 for supervised SimCSE,
76.25 for unsupervised SimCSE and 56.70 for BERT-base with first--last
averaging, all within 0.01 points of Table~5 in~\cite{gao-etal-2021-simcse}
(Section~\ref{sec:results});
\item re-implements the SimCSE loss independently and tests it against the loss
computation in the official code (Section~\ref{sec:method-infonce}).
% Number below: experiments/reproduction/results/R4-2026-11-24/metrics.json.
In float64 the two agree within $1.4\times10^{-9}$ over all tested batch
sizes, temperatures and hard-negative weights;
% Corpus: data/d2/documents.yaml (decision D-17 in docs/application/PROJECT.md).
\item specifies a retrieval assistant for the undergraduate training regulations of
our own university, the University of Science, VNU-HCM. Its corpus holds both
the 2021 regulation and the 2026 regulation that replaces it, because the 2026
text applies to cohorts admitted from the 2026--2027 academic year, while for
earlier cohorts it lets the university apply the new rules that are more
favourable to students. The assistant combines
BM25 with a SimCSE-style Vietnamese encoder, ranks whole articles, shows the
best-matching passage with its document, decision number and scope of
application, and flags low-confidence results instead of generating an answer
(Sections~\ref{sec:method-pipeline} and~\ref{sec:app});
% Checked 23/11 on ViRHE4QA_v2.csv @ ViRE 60ea715: 77 of the 297 distinct
% contexts name the university (ĐHCNTT / Trường ĐH CNTT); no other university
% is named (data/raw/virhe4qa/checks/outputs_2026-11-23.txt). The R2GQA abstract
% does not name the source university, so the attribution below rests on the
% article text.
\item defines an evaluation protocol on two datasets that the paper
does not use. The first is the public ViRHE4QA
subset~\cite{do2024r2gqaretrieverreadergeneratorquestionanswering,nguyen-quan-2026-works},
evaluated on its own corpus with a development/test split grouped by source
article. Judging by their wording, its articles are regulations of the
University of Information Technology (UIT), VNU-HCM, the only university that
they name. The second is a planned test set of student-style questions
about our university's two regulations, including questions that the
regulations do not answer. The protocol calls for paired bootstrap confidence
intervals. \TODO{Report the main retrieval finding once E1--E3 are complete.}
\end{itemize}

Section~\ref{sec:related} reviews related work. Section~\ref{sec:method}
describes SimCSE and our retrieval pipeline, and Section~\ref{sec:setup} gives
the datasets, metrics and protocols. Section~\ref{sec:results} reports the
reproduction results. The retrieval experiments and application evaluation
are pending.
